\documentclass{article}
\usepackage[margin=1in]{geometry}
\usepackage{longtable}
\usepackage{array}

\begin{document}

\section*{GeoTop Formalization Report for Sections 3 and 4}

Date: 2026-06-02.

Sections 3 and 4 of \texttt{geotop.tex} are not yet sorry-free.  They are
isolated in a cached development stack and the current active session builds
after the latest merge from \texttt{main}.  The merge brought in colleague
\texttt{main} at \texttt{3e463c3b}; the local merge commit is
\texttt{2b0745a0}.

The post-merge check passed:
\begin{itemize}
\item \texttt{/project/bin/isabelle build -d . -d dev34\_pre -d dev34\_prefix -d dev34\_facts -d dev34\_workfacts -d dev34\_linkfacts -d dev34\_graphfacts -d dev34\_graphwork -d dev34\_openstar -d dev34 GeoTop34Dev}
\end{itemize}

Because the merge touched base files, the full post-merge rebuild took
\texttt{0:11:32} elapsed.  The final \texttt{GeoTop34Dev} theory itself ran in
\texttt{0:00:07} once the caches had rebuilt.

Excluding the intentionally dirty \texttt{dev34\_pre/GeoTop.thy} mirror, the
target section-specific theories currently contain 18 executable
\texttt{sorry}s: 10 in \texttt{dev34\_prefix/GeoTop\_3\_4\_Prefix.thy} and 8 in
\texttt{dev34/GeoTop\_3\_4.thy}.

\section*{Layout}

The section 3--4 development is split as follows:
\begin{itemize}
\item \texttt{dev34\_pre/GeoTop.thy}: dirty cached import of the original script; it still contains original later-section sketches and is excluded from the Section 3--4 target count.
\item \texttt{dev34\_prefix/GeoTop\_3\_4\_Prefix.thy}: Section 3 and early Section 4.
\item \texttt{dev34\_facts/GeoTop\_3\_4\_Facts.thy}: Jordan-curve facts and local manifold helpers.
\item \texttt{dev34\_workfacts}, \texttt{dev34\_linkfacts}, \texttt{dev34\_graphfacts}, \texttt{dev34\_graphwork}, and \texttt{dev34\_openstar}: cached support layers.
\item \texttt{dev34/GeoTop\_3\_4.thy}: active Section 4 manifold/star work.
\end{itemize}

\section*{Section 3 Counterparts}

\begin{longtable}{>{\raggedright\arraybackslash}p{0.13\linewidth}
                  >{\raggedright\arraybackslash}p{0.45\linewidth}
                  >{\raggedright\arraybackslash}p{0.18\linewidth}
                  >{\raggedright\arraybackslash}p{0.16\linewidth}}
\textbf{Book item} & \textbf{Book statement, abbreviated} & \textbf{Isabelle} & \textbf{File}\\
\hline
Theorem 3.1 & Matched simplexes admit a simplicial homeomorphism carrying vertices to assigned images. & \texttt{Theorem\_GT\_3\_1} & \texttt{dev34\_prefix/GeoTop\_3\_4\_Prefix.thy}\\
Theorem 3.2 & In the equal-dimensional case, the simplex homeomorphism extends to the whole Euclidean space. & \texttt{Theorem\_GT\_3\_2} & \texttt{dev34\_prefix/GeoTop\_3\_4\_Prefix.thy}\\
Theorem 3.3 & A triangulated polygonal disk with more than one 2-simplex has a free 2-simplex. & \texttt{Theorem\_GT\_3\_3} & \texttt{dev34\_prefix/GeoTop\_3\_4\_Prefix.thy}\\
Theorem 3.4 & Every polygon is carried by a plane homeomorphism to the frontier of a 2-simplex. & \texttt{Theorem\_GT\_3\_4} & \texttt{dev34\_prefix/GeoTop\_3\_4\_Prefix.thy}\\
Theorem 3.5 & Any two polygons in \(R^2\) are equivalent by a plane homeomorphism. & \texttt{Theorem\_GT\_3\_5} & \texttt{dev34\_prefix/GeoTop\_3\_4\_Prefix.thy}\\
Theorem 3.6 & A polygon bounds a topological 2-cell. & \texttt{Theorem\_GT\_3\_6} & \texttt{dev34\_prefix/GeoTop\_3\_4\_Prefix.thy}\\
Theorem 3.7 & The polygon-to-simplex homeomorphism can be chosen fixed outside a prescribed open neighborhood. & \texttt{Theorem\_GT\_3\_7} & \texttt{dev34\_prefix/GeoTop\_3\_4\_Prefix.thy}\\
\end{longtable}

\section*{Section 4 Counterparts}

\begin{longtable}{>{\raggedright\arraybackslash}p{0.13\linewidth}
                  >{\raggedright\arraybackslash}p{0.45\linewidth}
                  >{\raggedright\arraybackslash}p{0.18\linewidth}
                  >{\raggedright\arraybackslash}p{0.16\linewidth}}
\textbf{Book item} & \textbf{Book statement, abbreviated} & \textbf{Isabelle} & \textbf{File}\\
\hline
Theorem 4.1 & Distinct components of an open Euclidean set are separated by disjoint open subsets. & \texttt{Theorem\_GT\_4\_1} & \texttt{dev34\_prefix/GeoTop\_3\_4\_Prefix.thy}\\
Theorem 4.2 & An arc across a polygonal disk separates the polygon interior into two expected open sides. & \texttt{Theorem\_GT\_4\_2} & \texttt{dev34\_prefix/GeoTop\_3\_4\_Prefix.thy}\\
Theorem 4.3 & A topological 1-sphere in \(R^2\) separates the plane. & \texttt{Theorem\_GT\_4\_3} & \texttt{dev34\_prefix/GeoTop\_3\_4\_Prefix.thy}\\
Theorem 4.4 & Two disjoint arcs from opposite boundary points leave \(Q\) and \(S\) in the frontier of the same component. & \texttt{Theorem\_GT\_4\_4} & \texttt{dev34\_prefix/GeoTop\_3\_4\_Prefix.thy}\\
Theorem 4.5 & No arc separates \(R^2\). & \texttt{Theorem\_GT\_4\_5} & \texttt{dev34\_facts/GeoTop\_3\_4\_Facts.thy}\\
Theorem 4.6 & For a 1-sphere \(J\) and component \(U\) of \(R^2-J\), \(\operatorname{Fr} U=J\). & \texttt{Theorem\_GT\_4\_6} & \texttt{dev34\_facts/GeoTop\_3\_4\_Facts.thy}\\
Theorem 4.7 & A 1-sphere in \(R^2\) has exactly one bounded complement component. & \texttt{Theorem\_GT\_4\_7} & \texttt{dev34\_facts/GeoTop\_3\_4\_Facts.thy}\\
Theorem 4.8 & If \(|K|\) is a 2-manifold, every vertex star of \(K\) is a combinatorial 2-cell. & \texttt{Theorem\_GT\_4\_8} & \texttt{dev34/GeoTop\_3\_4.thy}\\
Theorem 4.9 & If \(|K|\) is a 2-manifold with boundary, vertex stars are combinatorial 2-cells and the boundary is the union of one-triangle edges. & \texttt{Theorem\_GT\_4\_9} & \texttt{dev34/GeoTop\_3\_4.thy}\\
Theorem 4.10 & For a closed 2-manifold with boundary in \(R^2\), manifold boundary equals topological frontier. & \texttt{Theorem\_GT\_4\_10} & \texttt{dev34/GeoTop\_3\_4.thy}\\
\end{longtable}

\section*{Open Proof Holes}

The current holes in \texttt{dev34\_prefix/GeoTop\_3\_4\_Prefix.thy} are the
strong free-simplex induction for \texttt{Theorem\_GT\_3\_3}, the base and step
of the polygon-reduction induction for \texttt{Theorem\_GT\_3\_4}, the
support-in-\(U\) version \texttt{Theorem\_GT\_3\_7}, the final decomposition
step of \texttt{Theorem\_GT\_4\_2}, and five brick-decomposition/frontier steps
inside \texttt{Theorem\_GT\_4\_4}.

The current holes in \texttt{dev34/GeoTop\_3\_4.thy} are
\texttt{geotop\_punctured\_star\_radial\_endpoint\_projection\_dev34}, the two
punctured-star connectedness lemmas, the cone-over-link equivalence lemma,
three local chart contradiction lemmas for edge-incidence counts, and the
boundary equality half of \texttt{Theorem\_GT\_4\_9}.

\section*{Recent Progress and Supporting Material}

The active file now proves the radial part of Moise's Lemma 5 argument through
\texttt{geotop\_punctured\_star\_separation\_side\_meets\_vertex\_neighborhood\_dev34}.
Important supporting helpers include
\texttt{geotop\_complex\_vertex\_star\_neighborhood\_dev34},
\texttt{geotop\_star\_punctured\_point\_radial\_to\_link\_dev34},
\texttt{geotop\_link\_radial\_endpoint\_unique\_dev34},
\texttt{geotop\_plane\_chart\_open\_subset\_connected\_punctured\_neighborhood},
and \texttt{geotop\_2\_cell\_open\_subset\_connected\_punctured\_neighborhood}.

\section*{Notes For Future Work}

The next book-aligned step is to finish the two punctured-star connectedness
lemmas using the new side-accumulation bridge and the plane/2-cell punctured
neighborhood facts.  After that, finish connected links, then the cone-over-link
bridge, then the remaining edge-incidence chart contradictions and the boundary
equality in Theorem 4.9.

\end{document}
